\section{Experimental setup}
\label{sec:protocol}

Appendix~\ref{app:setup} gives every setting (Table~\ref{tab:setup}), arm
(Table~\ref{tab:arms}) and Holm family (Table~\ref{tab:families}), and the
recipes for CelebA-64, DreamBooth \citep{ruiz2023dreambooth} with Stable
Diffusion~1.5 \citep{rombach2022high} and denoising.

\subsection{Training}
\label{sec:setup-training}

All arms on CIFAR-10 \citep{krizhevsky2009learning} share the recipe of the
code released with the conference paper \citep{thao2026awwl} and, apart from
the two EMA arms, differ only in the objective. The denoiser is a DDPM-style
U-Net \citep{ho2020denoising}, the \texttt{UNet2DModel} of \texttt{diffusers},
with channels 128, 256, 256 and 256, two residual blocks per level, one
self-attention layer (in the $4\times4$ bottleneck only), no dropout and
34.4~million parameters. The process has $T = 1000$ steps with $\beta_t$
linear from $10^{-4}$ to $0.02$ and $\epsilon$-prediction
(Eq.~\eqref{eq:simple}). Training uses AdamW \citep{loshchilov2019decoupled}
at learning rate $2\times10^{-4}$ with 500 warm-up steps and cosine decay,
batch size 128 and float16 mixed precision, for 200 epochs (\num{78200}
steps). Exponential moving averaging (EMA, decay 0.9999)
\citep{polyak1992acceleration}, which the published runs omit, enters the
200-epoch matrix only as two separate arms. Training seeds 1 to 5 fix the
initialisation and the data order, and seed $k$ of an arm is paired with seed
$k$ of the arm it is compared with. All checkpoints are sampled from one
sampling seed, so every interval describes training variance conditional on
one draw of initial latents.

The arms (Table~\ref{tab:arms}) are MSE, three other pixel losses, two
timestep reweightings (Section~\ref{sec:notation}), FFL, AWWL with its static
($p = 0$) and summed-detail (Eq.~\eqref{eq:summed}) variants, four
normalisation controls, AWWL's timestep component alone (the timestep-only
arm; Section~\ref{sec:confound}), MSE and AWWL with EMA, and seven extensions
(Section~\ref{sec:extensions}). The EMA pair, the learned basis and the
all-axes combination have three seeds; the conference paper's VGG loss is
omitted (Appendix~\ref{app:setup-arms}). Three settings depart from the
matrix: 800 epochs with EMA (MSE and AWWL, three seeds), DDPM-1000 sampling
(MSE, the static wavelet arm and AWWL) and CelebA \citep{liu2015deep} at
$64\times64$ (MSE, five seeds).

\subsection{Sampling and metrics}
\label{sec:setup-metrics}

Each evaluated checkpoint generates \num{10000} samples with DDIM
\citep{song2021denoising} (100 steps, $\eta = 0$), stored as 8-bit PNG. The
main operating point is epoch 199, the last of the 200 epochs counted from
zero; MSE, the static wavelet arm and AWWL are also evaluated at epochs 39,
79, 119 and 159. FID \citep{heusel2017gans} is computed with \texttt{clean-fid}
\citep{parmar2022aliased} in \texttt{clean} mode against all
\num{\nRefImages} training images. FID from \num{10000} samples is biased
upward relative to FID-50k by a generator-dependent amount
\citep{chong2020effectively}, which need not cancel between arms, and
absolute values are not comparable with FID-50k values in the literature.
Inception Score \citep{salimans2016improved}, Kernel Inception Distance
(KID) \citep{binkowski2018demystifying}, precision and recall
\citep{kynkaanniemi2019improved} are reported alongside FID
(Table~\ref{tab:full}; Appendix~\ref{app:setup-metrics}), and the pricing
test is repeated in DINOv2 features (Section~\ref{sec:worth-dino}). Operator
outputs of the pricing test are scored in floating point, without rounding to
8 bits (hazard H1, Section~\ref{sec:worth-curve}). The
spectral statistic $s$ (Eq.~\eqref{eq:statistic}) is computed on all samples
and reference images, not on the first files of a directory listing, a
subset that is neither random nor portable (hazard H2,
Section~\ref{sec:mechanism}).

\subsection{Statistical analysis}
\label{sec:setup-stats}

Comparisons report the mean paired difference over seeds with its 95\%
Student-$t$ interval and the two-sided paired $t$-test, which is exact
whatever the correlation between arms (Appendix~\ref{app:setup-stats});
``$\pm$'' is the sample standard deviation across seeds. The Wilcoxon
signed-rank test \citep{wilcoxon1945individual} checks the direction only:
its smallest two-sided $p$, $2/2^{n}$, is $0.0625$ for five pairs and $0.25$
for three, so any smaller $p$, including every significant result at three
seeds, rests on the $t$ model. Holm--Bonferroni correction
\citep{holm1979simple} is applied within fixed families, one per question
(Table~\ref{tab:families}), and raw and corrected $p$ are both reported. Each
hypothesis belongs to one family: the budget family holds only the two
800-epoch contrasts, and the timestep-only arm has a two-test family of its
own (the arm against MSE, and AWWL against the arm). The families are fixed
in the analysis code, so a corrected $p$ does not depend on which runs have
finished; they were set during the analysis, not before it. AWWL's MRG
(Eq.~\eqref{eq:mrg}), the headline quantity, is corrected over the
\nFamilyPricing{} arms of the pricing family, and its FID gap to MSE over the
\nFamilyMain{} arms of the main family.
Convergence, band-wise, operator and positive-control results carry
intervals and uncorrected $p$ (for the null control, in the pricing
family, also Holm-adjusted $p$). An unresolved (non-significant) difference is read
through the effects its interval excludes; at five seeds and the spread of the
AWWL$-$MSE differences, a paired test has 80\% power against a true FID
difference of about \mdeAWWL{} (Appendix~\ref{app:setup-stats}).

Every table and measured number is generated by scripts from the run ledgers,
released with the code (Data availability); an AI coding assistant wrote the
code under the authors' supervision.
